\documentclass[10pt,a4paper]{article}
\usepackage{amsmath,amssymb,amsthm,graphicx}
% Minimal retained preamble for lane-p2 cards taken from
% sources/ahl-16-fixed-points-2019/editorial-preamble.tex
% (doi:10.5802/ahl.16, Annales Henri Lebesgue 2 (2019) 149-185, CC BY 4.0).
%
% Copied verbatim from the editorial preamble: the package list, the
% theorem declarations under the shared counter and the macro block, so that
% every macro and environment of the retained spans is available.
%
% Omitted and why: the class line (the assembler supplies the class); the
% a4paper geometry line and the fontspec block with its explicit lmroman10
% files (encodings and fonts fight the forced xelatex engine, and the
% assembler supplies the page setup).
\usepackage{amsmath,amssymb,amsthm,mathtools,graphicx,xcolor,hyperref,xurl,xspace,ifthen}
\newtheorem{theo}{Theorem}[section]
\newtheorem{lemm}[theo]{Lemma}
\newtheorem{theorem}[theo]{Theorem}
\newtheorem*{remas*}{Remarks}
\providecommand{\og}{«\nobreakspace}
\providecommand{\fg}{\nobreakspace»}
\newtheorem{lem}[theo]{Lemma}
\newtheorem{lemma}[theo]{Lemma}
\newtheorem{prop}[theo]{Proposition}
\newtheorem{coro}[theo]{Corollary}
\newtheorem{corol}[theo]{Corollary}
\newtheorem{conj}[theo]{Conjecture}
\newtheorem{Proposition}[theo]{Proposition}
\newtheorem{theoreme}[theo]{Theorem}
\theoremstyle{definition}
\newtheorem{defi}[theo]{Definition}
\newtheorem{exam}[theo]{Example}
\newtheorem{exem}[theo]{Example}
\newtheorem{rema}[theo]{Remark}
\newtheorem{nota}[theo]{Notation}
\newtheorem*{rema*}{Remark}
\newtheorem*{defi*}{Definition}
\newtheorem*{theo*}{Theorem}
\newtheorem*{lemm*}{Lemma}
\newtheorem*{prop*}{Proposition}
\newcommand{\enoncename}{}
\newtheorem{corpusenonce}[theo]{\enoncename}
\newenvironment{enonce}[2][]{\renewcommand{\enoncename}{#2}\begin{corpusenonce}}{\end{corpusenonce}}
\newenvironment{enonce*}[2][]{\par\medskip\noindent\textbf{#2.}\itshape}{\par\medskip}
\newenvironment{proofc}[1][\proofname]{\begin{proof}[#1]}{\end{proof}}
\newcommand{\Mc}{\mathcal{M}}
\newcommand{\pointir}{.\ ---\ }
\providecommand{\mathof}[1]{\mathrm{#1}}

\numberwithin{equation}{section}
\numberwithin{figure}{section}
\numberwithin{table}{section}
\def\endappendix{}
%~Mouliné par MaN_auto v.0.21.4 2019-04-30 09:02:09


\usepackage[all]{xy}

\hyphenation{auto-morphism
every
Roelcke
peri-od-ic
peri-od-icy
rep-re-sen-ta-tion
}








\newcommand{\A}{\mathcal A}
\newcommand{\N}{\mathbb{N}}
\newcommand{\Z}{\mathbb{Z}}
\newcommand{\Q}{\mathbb{Q}}
\newcommand{\K}{\mathcal{K}}
\newcommand{\CC}{\mathbb{C}}

\newcommand*{\mA}{\mathbf{A}}
\newcommand*{\mB}{\mathbf{B}}
\newcommand*{\mC}{\mathbf{C}}
\newcommand*{\mD}{\mathbf{D}}
\newcommand*{\mE}{\mathbf{E}}
\newcommand*{\mF}{\mathbf{F}}
\newcommand*{\mH}{\mathbf{H}}
\newcommand*{\mZ}{\mathbf{Z}}
%\newcommand*{\mA}{\mathbf{A}}
%\newcommand*{\mA}{\mathbf{A}}

\newcommand{\restrict}[2]{#1 \upharpoonright #2}

\newcommand{\e}{\varepsilon}

\DeclareMathOperator{\Aut}{Aut}
\DeclareMathOperator{\Stab}{Stab}
\DeclareMathOperator{\Age}{Age}

\newcommand{\actson}{\curvearrowright}

\newcommand{\RUCB}{\mathrm{RUC}_{b}}
\newcommand{\ROB}{\mathrm{Ro}_{b}}
\newcommand{\WAP}{\mathrm{WAP}}
\newcommand{\AP}{\mathrm{AP}}

\graphicspath{{./figures/}}

\newcommand*{\mk}{\mkern -1mu}
\newcommand*{\Mk}{\mkern -2mu}
\newcommand*{\mK}{\mkern 1mu}
\newcommand*{\MK}{\mkern 2mu}

\hypersetup{urlcolor=purple, linkcolor=blue, citecolor=red}

\newcommand*{\romanenumi}{\renewcommand*{\theenumi}{\roman{enumi}}}
\newcommand*{\Romanenumi}{\renewcommand*{\theenumi}{\Roman{enumi}}}
\newcommand*{\alphenumi}{\renewcommand*{\theenumi}{\alph{enumi}}}
\newcommand*{\Alphenumi}{\renewcommand*{\theenumi}{\Alph{enumi}}}
\let\oldexists\exists
\renewcommand*{\exists}{\mathrel{\oldexists}}
\let\oldforall\forall
\renewcommand*{\forall}{\mathrel{\oldforall}}


\usepackage{cleveref}
\crefname{originalproposition}{original proposition}{original propositions}
\Crefname{originalproposition}{Original proposition}{Original propositions}
\makeatletter
\newlabel{prop:Gelfand}{{3.9}{source}{}{sourcepointer-1}{}}
\newlabel{prop:Gelfand@cref}{{[originalproposition][0][]3.9}{[1][0][]source}{}{}{}}
\makeatother
\title{A class of flows closed under suprema and factors yields a left-invariant algebra whose factors are exactly its ambits}
\author{Lionel Nguyen Van Thé}
\date{doi:10.5802/ahl.16; CC BY 4.0}
\begin{document}
\maketitle
\noindent\textit{Editorial note.} The counted unit is the article's own Proposition prop:KA at source lines 457-459 together with its proof at source lines 461-463, a straight adjacent pairing in the source file. The source is the publisher version of the article in Annales Henri Lebesgue 2 (2019), pages 149-185, DOI 10.5802/ahl.16, released under CC BY 4.0.\par\medskip
\noindent\textit{Editorial note.} The retained proof invokes the article's Gelfand correspondence. That statement is not retained here, so it is declared as a navigation pointer with the printed number 3.9, read off an enumeration of every numbered theorem-like environment of the source in order. It is a statement this lane has already counted on another card of the same source; the pointer still carries only its number and none of its text, which is the same treatment every other pointer of this lane receives.\par\medskip
\noindent\textit{Editorial note.} The licence evidence used here is the observation file kept with the source, license-observation.json, which records the licence as CC BY 4.0, points at the inherited publisher-version evidence and states that the landing page displays no licence line; the underlying evidence files are the same ones the earlier cards of this source rely on. No acceptance record is written anywhere.\par\medskip
\noindent\textit{Editorial note.} The retained preamble is this source's editorial preamble with three omissions only: the class line, which the assembler supplies; the page geometry; and the fontspec block with its explicit font files, which would fight the forced xelatex engine. Everything else, including the theorem declarations under the shared counter and the whole macro block, is retained verbatim. This card copies the shared, card-independent slice kept at focused-seeds/logic/P2-ahl16-retained-preamble.tex, which is the same 123 lines every other card of this source uses.\par\medskip
\noindent\textit{Editorial note.} The gate raised one manual-review note: the next proof environment begins at line 465. That block belongs to the following statement of the section and is not part of this card; the proof retained here is closed at line 463.\par\medskip
\noindent\textit{Editorial note.} Neither the retained statement nor the retained proof contains a citation, so this card needs no bibliography entry.\par\medskip
\noindent\textit{Editorial note.} The census package marked this candidate free of missing citation keys and listed the statement span as 457-459 and the proof span as 461-463; both were re-read from the source before the card was written.\par\medskip
\noindent\textit{Editorial note.} rqual sv-mismatch batch fix 2026-10-14 (Chief ruling: re-assemble unpublished only): source\_version aligned verbatim to the source dir license-observation.json sv (package evidence authority); revision bumped for re-assembly. Recipe from lane-s3.\par\medskip
\noindent\textit{Allowed standard prerequisites (editorial).}\begin{itemize}
\item G-flows, ambits, suprema and factors. An ambit is a flow with a distinguished point whose orbit is dense, the class of ambits of a class of flows is its closure under these maps, suprema and factors are the two operations of the category of flows, and a closed left-invariant unital star-subalgebra of the bounded right-uniformly continuous functions carries the flow of the group on its spectrum, whose ambits are the factors of that flow. The author's own definitions are the ones used by the counted statement and proof.. The counted proposition is the construction of the algebra attached to a class of flows, and it states that the class of flows closed under suprema and factors is recovered from the algebra as the family of factors of one canonical flow.
\end{itemize}\medskip
\ifdefined\hypertarget\hypertarget{sourcepointer-1}{}\fi
\noindent\textit{Source reference.} Original proposition 3.9: sources/ahl-16-fixed-points-2019/body-01.tex, source lines 421--423. Navigation pointer invoked by the retained proof at the step that passes to Gelfand compactifications: the containment of ambits is turned into being a factor of the canonical flow there. That proposition is an earlier result of the article and is not retained in this card, so the reference is declared as a pointer carrying the number the article itself prints for that statement, that is the section and the place of the statement inside that section, counted by enumerating every numbered theorem-like environment of the source in order: it is 3.9; the cited statement is an allowed input and none of it is reproduced inside the span.\par\medskip
\begin{prop}\label{prop:KA}
Let $G$ be a topological group and $\mathcal X$ be a class of $G$-flows such that the class of $\mathcal X$-$G$-ambits is closed under suprema and factors. Then the set $\A=\{ f\in \RUCB(G) : G\actson \overline{G\bullet f} \in \mathcal X\}$ forms a unital left-invariant, closed C$^{*}$-subalgebra of $\RUCB(G)$, and the factors of $G\actson(G^{\A}, e_{G})$ are exactly the $\mathcal X$-$G$-ambits.
\end{prop}

\begin{proof}
Let $(X,x)=\bigvee_{f\in \A}(\overline{G\bullet f},f)$. As a supremum of $\mathcal X$-$G$-ambits, it is a $\mathcal X$-$G$-ambit as well. Let $C(X)_{x} = \{ f_{x}:f\in C(X)\}$. As we have seen, this is a unital left-invariant, closed C$^{*}$-subalgebra of $\RUCB(G)$. To prove the result, it suffices to show that it is equal to $\A$. Let $f\in \RUCB(G)$. Then $f\in C(X)_{x}$ iff $\langle f\rangle \subset C(X)_{x}$. Passing to Gelfand compactifications, this means that $(G^{\langle f \rangle},e_{G})$ is a factor of $(G^{C(X)_{x}},e_{G})$, or, equivalently, that $(\overline{G\bullet f},f)$ is a factor of $(X,x)$ (Proposition~\ref{prop:Gelfand}). Now, this happens iff $f\in \A$: the direct implication holds because the class of $\mathcal X$-$G$-ambits is closed under factors, and the converse holds thanks to the definition of $(X,x)$, as $(\overline{G\bullet f},f)$ appears as one of its factors.
\end{proof}

\end{document}
